% 05. kapitulua - Ariketak: zereginak eta taulak GitHuben
% Helburua: agent-scrumera iritsi aurretik eskuz praktikatzea agenteek gero kontsumituko duten guztia:
%           GitHub Projects taula bat sortu, Kanban eta Scrum ikuspegiak gehitu, issue-ak idatzi eta
%           Egiteko / Egiten / Eginda zutabeetan zehar mugitu, adarra eta pull request-a issue bati lotu.
%           Kapitulu praktikoa, teoria berririk gabe: 1-4. kapituluetako kontzeptuak eskuekin finkatzeko.
% Ariketa-sekuentzia proposatua (aurreko liburuko formatuan: "N. ariketa - ...", "Ariketa ebatziak", "Proposatutako ariketa gehigarriak"):
%   1. ariketa - Taldearen biltegia eta GitHub Projects taula sortu, biltegiari lotu, taldekideak gehitu.
%   2. ariketa - Kanban ikuspegia: Egiteko / Egiten / Eginda zutabeak (Status eremua), "Egiten" zutabeari muga jarri (WIP).
%   3. ariketa - Scrum ikuspegia: Iteration (sprint) eremua, Priority eta Estimate eremu pertsonalizatuak,
%                taula-ikuspegia iterazioka taldekatuta, Roadmap ikuspegia.
%   4. ariketa - Issue-ak idatzi: issue-txantiloia (testuingurua, onarpen-irizpideak, "eginda" definizioa),
%                etiketak (bug, feature, docs), arduraduna, mugarria (milestone); ekosistemaren backlog-etik abiatuta.
%   5. ariketa - Issue-ak eskuz mugitu zutabeetan zehar; egoera-aldaketa bakoitza iruzkin batekin dokumentatu (egoera-txostena).
%   6. ariketa - Automatizazioak: Projects-en workflow integratuak (issue berria -> Egiteko; issue itxita edo PR bateratuta -> Eginda).
%   7. ariketa - Issue batetik adarra sortu, konfirmazioetan issue-a aipatu (#N), pull request-a "Closes #N"-rekin,
%                berrikuspena, bateratzea; egiaztatu issue-a itxi dela eta txartela Eginda zutabera mugitu dela.
%   8. ariketa - Sprint baten simulazioa rol-jokoan: Nami, Ane, Mikel eta Gotzon; sprint backlog-a, eguneroko taula-bilera,
%                burndown (Insights), sprint-itxiera eta atzera-begirakoa. Gizakiek egiten dute gero agenteek egingo dutena.
%   9. ariketa - Blokeoak eta eskalatzea: "blocked" eta "needs-human" etiketak, issue-aren egoera-makina paperean
%                (7. kapituluan formalizatuko dena).
%  10. ariketa - Gauza bera terminaletik: gh CLI (gh issue create, gh issue edit, gh project item-add, gh pr create).
%                Hau da agenteek gero erabiliko duten interfazea: zubia 6. kapitulura.
% Aurreko liburutik berrerabil daitekeena:
%   - aurreko_liburua_2024/kapituluak/9-ariketa_adarkatze.tex (product/sprint backlog, taldea, zikloak)
%   - aurreko_liburua_2024/kapituluak/5-prestaketa_lanak.tex (ariketen formatua)
% Irudiak: Pictures/05/ (goiburu hutsa header_05.pdf.xoj); backlog-aren irudiak Pictures/09/ karpetan daude (9.backlog, 9.product_backlog.xoj)
\todo[inline]{Idazteko.}
